% ════════════════════════════════════════════════════════════════════
% Chapter 2 — System Architecture
% ════════════════════════════════════════════════════════════════════
\section{System Architecture}\label{sec:arch}

\subsection{Overview}

The data collection platform is designed around a physically distributed, wireless topology that fully decouples the barbell-mounted sensing hardware from the host workstation.
This decoupling is a strict biomechanical requirement: any physical cable attached to the barbell introduces spurious restoring forces, alters the natural bar path, and compromises the ecological validity of the kinematic measurements.
The system comprises four distinct functional units, depicted schematically in \cref{fig:arch}.

\begin{figure}[H]
  \centering
  \resizebox{\textwidth}{!}{
  \begin{tikzpicture}[
    node distance=2.5cm,
    >={Stealth[length=3mm]},
    block/.style={draw, rectangle, minimum width=4cm, minimum height=2.5cm, align=center, rounded corners=2pt, font=\small},
    sensor/.style={draw, rectangle, minimum width=3cm, minimum height=1.8cm, align=center, rounded corners=2pt, font=\small, fill=blue!5},
    wire/.style={->, thick},
    wireless/.style={->, thick, decorate, decoration={snake, amplitude=1.5pt, segment length=8pt}, color=purple!70!black},
    sync/.style={->, thick, dashed, color=red!70!black},
    label/.style={font=\scriptsize, align=center},
    ]

    % === Camera ===
    \node[block, fill=orange!10] (cam) {
      \textbf{Intel RealSense D455}\\[2pt]
      Depth Camera\\
      $848\times480$ @ \SI{90}{fps}\\
      \textit{Timing Master}
    };

    % === Barbell Node ===
    \node[block, fill=green!8, right=4cm of cam] (esp_bar) {
      \textbf{Barbell Node}\\[2pt]
      ESP32-WROOM-32\\
      160\,MHz, LiPo powered\\
      BOD disabled
    };

    % === IMU ===
    \node[sensor, below=2cm of esp_bar] (imu) {
      \textbf{ICM-42688-P}\\[2pt]
      6-axis IMU\\
      \SI{1}{\kilo\hertz} ODR
    };

    % === Level Shifter ===
    \node[draw, diamond, aspect=2, fill=yellow!15, font=\scriptsize, below=1.8cm of cam, align=center] (ls) {
      Level Shifter\\
      2$\times$2N3904
    };

    % === Relay ===
    \node[block, fill=purple!8, right=4cm of esp_bar] (relay) {
      \textbf{Relay Node}\\[2pt]
      ESP32-WROOM-32\\
      USB-connected\\
      Ring-buffered
    };

    % === Host PC ===
    \node[block, fill=gray!10, right=3.5cm of relay] (pc) {
      \textbf{Host PC}\\[2pt]
      C++17 Application\\
      Python Golden Model\\
      Data Storage
    };

    % === Wires ===

    % Camera -> Level Shifter
    \draw[sync] (cam.south) -- node[left, label] {\SI{1.8}{\volt} SYNC\\Pin 5 Aux} (ls.north);

    % Level Shifter -> IMU FSYNC
    \draw[sync] (ls.east) -| node[near end, right, label] {\SI{3.3}{\volt}\\FSYNC} (imu.west);

    % Level Shifter -> ESP32 GPIO
    \draw[sync] (ls.north east) -- node[above, label, sloped] {\SI{3.3}{\volt} GPIO\,27} (esp_bar.south west);

    % ESP32 <-> IMU (SPI)
    \draw[wire, <->] (esp_bar.south) -- node[right, label] {SPI @ \SI{8}{\mega\hertz}\\14-byte burst} (imu.north);

    % ESP32 -> Relay (ESP-NOW)
    \draw[wireless] (esp_bar.east) -- node[above, label] {ESP-NOW unicast\\208\,B @ \SI{125}{\hertz}} (relay.west);

    % Relay -> PC (USB)
    \draw[wire] (relay.east) -- node[above, label] {USB UART\\921\,600 baud} (pc.west);

    % Camera -> PC (USB 3.0)
    \draw[wire] (cam.east) -- ++(1,0) |- node[near start, above, label] {USB 3.0 Type-C} ([yshift=1.2cm]pc.west) -- (pc.north west);

    % D455 onboard IMU
    \node[sensor, fill=orange!5, below=2cm of cam, xshift=-3.5cm, minimum width=2.5cm, minimum height=1.3cm] (cam_imu) {
      \textbf{D455 BMI085}\\
      Onboard IMU
    };
    \draw[wire, dotted] (cam.south west) -- node[left, label] {Internal} (cam_imu.north);
    \draw[wire, dotted] (cam_imu.east) -- ++(1.5,0) |- node[near end, below, label] {camera\_imu.csv} ([yshift=-0.5cm]pc.south west);

    % Bounding box for barbell assembly
    \begin{scope}[on background layer]
      \node[draw=green!50!black, dashed, thick, rounded corners, fit=(esp_bar)(imu), inner sep=0.5cm, label={[font=\footnotesize\itshape]above:Barbell Assembly (battery-powered, wireless)}] {};
    \end{scope}

  \end{tikzpicture}
  }
  \caption{Complete system architecture showing data buses, wireless telemetry, and hardware synchronisation paths. Red dashed lines indicate the \SI{1.8}{\volt}$\rightarrow$\SI{3.3}{\volt} level-shifted SYNC pulse. Purple wavy lines indicate wireless ESP-NOW transmission.}
  \label{fig:arch}
\end{figure}


\subsection{Component Selection Rationale}

\subsubsection{Intel RealSense D455 Stereo Depth Camera}
The D455 was selected as the optical ground-truth sensor for the following reasons:
\begin{enumerate}[nosep]
  \item \textbf{Global shutter:} Both stereo IR imagers use global shutters, eliminating the rolling-shutter distortion that corrupts velocity estimation during fast barbell transients. At a peak velocity of \SI{2}{\metre\per\second}, a rolling shutter with a \SI{10}{\milli\second} readout produces up to \SI{20}{\milli\metre} of intra-frame geometric distortion.
  \item \textbf{Hardware sync output:} The D455 exposes an auxiliary connector with a dedicated \texttt{FRAME\_SYNC} output (Pin~5), enabling it to serve as the absolute timing master for the distributed system. This output is not available on consumer-grade depth cameras (e.g., Azure Kinect, Intel D435).
  \item \textbf{Wide stereo baseline:} The \SI{95}{\milli\metre} baseline between the left and right IR imagers provides reliable depth estimation up to \SI{6}{\metre}, easily covering the operating volume of a standard power rack (\SI{2}{}--\SI{3}{\metre} camera-to-barbell distance).
  \item \textbf{Onboard IMU:} The D455 integrates a Bosch BMI085 6-axis IMU, providing accelerometer data at \SI{250}{\hertz} and gyroscope data at \SI{200}{\hertz}. While this IMU is not used for VBT kinematics (the camera body is stationary on a tripod), it provides free cross-stream synchronisation validation and tripod-shake/bump detection.
  \item \textbf{Cost:} Approximately \$350~USD, three orders of magnitude cheaper than a multi-camera optoelectronic installation.
\end{enumerate}

\subsubsection{ESP32-WROOM-32 Microcontroller}
The Espressif ESP32 was chosen over alternative microcontrollers (e.g., STM32, nRF52, RP2040) specifically for its \textbf{ESP-NOW} protocol:
\begin{itemize}[nosep]
  \item ESP-NOW operates directly on the IEEE 802.11 MAC layer, bypassing the TCP/IP and WiFi association overhead entirely. This yields sub-\SI{2}{\milli\second} one-way latency with acknowledgement and automatic retry.
  \item Traditional BLE (Bluetooth Low Energy) is fundamentally unsuitable for \SI{1}{\kilo\hertz} streaming: BLE connection intervals are negotiated in multiples of \SI{1.25}{\milli\second} with typical minimums of \SI{7.5}{\milli\second}, creating severe bottlenecks. Even with BLE 5.0's extended data length, sustained throughput above \SI{20}{\kilo\byte\per\second} requires careful negotiation and is not reliable under the electromagnetic interference present in a gym environment.
  \item ESP-NOW supports unicast with hardware-level ACK, achieving $>99\%$ packet delivery at maximum TX power (20\,dBm).
\end{itemize}

\subsubsection{TDK InvenSense ICM-42688-P IMU}
The ICM-42688-P is a premium 6-axis MEMS MotionTracking device. It was selected for:
\begin{itemize}[nosep]
  \item Ultra-low gyroscope noise density: \SI{2.8}{\milli\degree\per\second}/$\sqrt{\text{Hz}}$, which is critical for accurate attitude estimation over multi-second integration windows.
  \item Configurable digital anti-aliasing filters: 1st, 2nd, or 3rd order, with bandwidth selectable as a fraction of the ODR.
  \item Internal FSYNC tagging engine: external synchronisation pulses are automatically tagged into the data payload (TEMP register LSB), eliminating the need for additional host-side interrupt handling.
  \item Wide full-scale ranges: $\pm\SI{16}{\g}$ accelerometer and $\pm\SI{2000}{\degree\per\second}$ gyroscope, sufficient for even the most explosive Olympic weightlifting movements.
\end{itemize}


\subsection{Data Flow}

The data flow through the system proceeds as follows:
\begin{enumerate}[nosep]
  \item The D455 camera generates a \texttt{FRAME\_SYNC} pulse every \SI{11.111}{\milli\second} (\SI{90}{\hertz}).
  \item This pulse passes through the 2$\times$2N3904 level shifter, producing a clean \SI{3.3}{\volt} rising edge.
  \item The shifted pulse simultaneously triggers: (a)~a rising-edge ISR on ESP32 GPIO~27, which records a 64-bit \si{\micro\second} timestamp; and (b)~the ICM-42688-P's internal FSYNC latch, which sets the LSB of the temperature register in the next data-ready sample.
  \item At \SI{1}{\kilo\hertz}, the ICM-42688-P asserts \texttt{INT1} upon each new sample. The ESP32 ISR sets a boolean flag; the main loop then executes a 14-byte SPI burst read, packs the data into a 26-byte CRC-protected packet, and appends it to a batching buffer.
  \item Every 8~samples (\SI{8}{\milli\second}), the 208-byte batch is transmitted via ESP-NOW unicast to the relay node.
  \item The relay node receives the ESP-NOW packet in an ISR, appends it to an 8\,KB ring buffer, and the main loop drains the buffer to the USB UART at 921\,600~baud.
  \item The host PC ingests both the serial IMU stream and the USB~3.0 camera stream. The \class{SyncEngine} class maintains a live affine mapping between ESP timestamps and camera hardware timestamps using the FSYNC tag events.
\end{enumerate}


\subsection{Latency Budget}

\Cref{tab:latency} summarises the end-to-end latency from physical sensor event to host-side availability.

\begin{table}[H]
  \centering
  \caption{End-to-end latency budget from sensor event to host availability.}
  \label{tab:latency}
  \begin{tabular}{llr}
    \toprule
    \textbf{Stage} & \textbf{Source} & \textbf{Latency} \\
    \midrule
    IMU ADC conversion         & ICM-42688-P internal      & $<\SI{100}{\micro\second}$ \\
    SPI burst read (14\,B)     & ESP32 @ \SI{8}{\mega\hertz}  & $\sim\SI{14}{\micro\second}$ \\
    Packet assembly + CRC      & ESP32 CPU                  & $\sim\SI{5}{\micro\second}$ \\
    Batch accumulation (8 pkt) & ESP32 RAM                  & \SI{8}{\milli\second} (worst case) \\
    ESP-NOW TX + ACK           & 802.11 MAC                 & $\sim\SI{2}{\milli\second}$ \\
    Relay ring buffer drain    & Serial @ 921\,600 baud     & $\sim\SI{2.3}{\milli\second}$ \\
    USB/OS scheduling jitter   & Host OS                    & $\sim\SI{1}{\milli\second}$ \\
    \midrule
    \textbf{Total (worst case)} & & $\sim$\SI{13.4}{\milli\second} \\
    \bottomrule
  \end{tabular}
\end{table}

Critically, the latency between sensor event and host availability is \emph{irrelevant} for post-processing accuracy, because every sample carries its own 64-bit ESP32 hardware timestamp (\code{esp\_timer\_get\_time()}), captured at the moment of the SPI read---before any batching or transmission.
The latency budget matters only for the real-time GUI feedback loop (operator view, live velocity plot), where \SI{13}{\milli\second} is imperceptible.
